\documentclass[11pt,a4paper]{article}
\usepackage[utf8]{inputenc}
\usepackage[T1]{fontenc}
\usepackage[portuguese]{babel}
\usepackage{lmodern}
\usepackage[margin=2.3cm]{geometry}
\usepackage{tabularx,booktabs,array,enumitem}
\usepackage[table]{xcolor}
\usepackage{textcomp}
\usepackage{hyperref}

\definecolor{roamly}{HTML}{0E6E68}
\definecolor{pendente}{HTML}{6B7280}
\hypersetup{colorlinks=true,linkcolor=roamly,urlcolor=roamly}
\newcolumntype{L}{>{\raggedright\arraybackslash}X}
\renewcommand{\arraystretch}{1.3}
\setlist{nosep,leftmargin=*}
\setlength{\parindent}{0pt}
\setlength{\parskip}{4pt}

% reference to an interview record or a breakdown (as in the report)
\newcommand{\reg}[1]{{\footnotesize\textcolor{pendente}{(#1)}}}
% use case
\newenvironment{usecase}[2]{%
  \par\medskip\noindent{\large\bfseries\color{roamly}#1: #2}\par\smallskip\small\noindent\ignorespaces}{\par\medskip}
\newcommand{\ucField}[2]{\textbf{#1} & #2 \\}

\begin{document}

\begin{usecase}{UC-04}{Levar o grupo a decidir sobre a proposta}
\textbf{Descrição breve.} O organizador do grupo partilha a proposta de viagem com os amigos que vão. Cada membro indica em que datas pode ir e se aceita cada item, e o organizador regista o que o grupo combinou. O sistema mostra a todos a mesma proposta e o mesmo total por pessoa, mostra ao organizador quem já respondeu e onde o grupo diverge, e guarda a decisão na proposta. A seguir, o organizador pode sair do Roamly para reservar no site de cada fornecedor.\par\smallskip
\begin{tabularx}{\textwidth}{@{}>{\raggedright\arraybackslash}p{3.4cm}L@{}}
\toprule
\ucField{Ator principal}{Organizador(a) do grupo}
\ucField{Atores secundários}{Membros do grupo (indicam as suas datas e a sua opinião sobre cada item)}
\ucField{Nível}{Resumo: a decisão demora dias e faz-se em várias sessões, de pessoas diferentes \reg{UE08}}
\ucField{Âmbito}{Roamly (versão base, viajante)}
\ucField{Área de foco}{5. Decidir com o grupo ou o cliente}
\ucField{Gatilho}{A proposta cabe no limite por pessoa (UC-02) e o organizador quer a decisão do grupo antes de o preço mudar.}
\ucField{Pré-condições}{A viagem tem uma proposta com o total por pessoa (UC-02) e os seus limites (UC-01). O organizador sabe quem vai.}
\ucField{Garantia mínima}{As datas e as opiniões já dadas ficam guardadas. Todos os membros veem a mesma versão, atual, da proposta.}
\ucField{Garantia de sucesso}{O que o grupo combinou fica registado na proposta, junto dos itens a que se refere, e todos os membros o veem. O organizador sabe a resposta de cada um e pode sair do Roamly para reservar no site de cada fornecedor.}
\bottomrule
\end{tabularx}

\medskip
\textbf{Partes interessadas e interesses}
\begin{itemize}
  \item \textbf{Organizador(a) do grupo:} ter uma decisão antes de o preço subir \reg{UE08}; saber a quem insistir sem perguntar um a um \reg{UE08}; não ser o único a lembrar-se, depois de cada conversa, do que ficou combinado \reg{UE18, U0-19}.
  \item \textbf{Membros do grupo:} responder sem abrir folhas de cálculo («manda print») \reg{UE17}; dizer o que não aceitam («dormitório não») antes de se reservar alguma coisa \reg{UE05}; dizer quando podem ir \reg{UE02, UE09}.
\end{itemize}

\noindent\begin{minipage}{\textwidth}
\setlength{\parskip}{4pt}
\textbf{Cenário principal de sucesso}

Cada passo está na coluna de quem o faz. Os passos 3 e 4 são feitos por cada membro, no seu tempo.

\begin{tabularx}{\textwidth}{@{}>{\raggedleft\arraybackslash}p{0.5cm}LLL@{}}
\toprule
 & \textbf{Organizador(a) do grupo} & \textbf{Membros do grupo} & \textbf{Sistema} \\
\midrule
1. & Partilha a proposta com o grupo. & & \\
2. & & & Mostra a proposta a cada membro: itens, preços, total por pessoa e o que foi cortado e porquê. \\
3. & & Cada membro indica se pode ir nas datas propostas ou em que outras datas pode. & \\
4. & & Cada membro indica se aceita cada item. & \\
5. & & & Mostra ao organizador quem já respondeu, quem ainda falta e os itens em que os membros divergem. \\
6. & Regista na proposta o que o grupo combinou. & & \\
7. & & & Guarda a decisão na proposta e mostra-a a todos os membros, junto dos itens a que se refere. \\
\bottomrule
\end{tabularx}
\end{minipage}

\textbf{Extensões}
\begin{itemize}[leftmargin=2.2em]
  \item[2a.] \textit{Um membro não abre a proposta partilhada} («manda print», \reg{UE17; FM.9, SM.21, AM.25}): o organizador envia um resumo com o total por pessoa que se lê na própria conversa do grupo (US-14). Continua no passo 5, onde esse membro continua a aparecer em falta.
  \item[3a.] \textit{Um membro só pode ir noutras datas} (o Amigo~B só pode a partir de uma sexta-feira concreta, \reg{UE09, U0-03}): o sistema regista essas datas como alternativa às propostas (US-15) e mostra-as ao organizador junto da resposta desse membro (US-31). Se o organizador refizer a pesquisa, a janela de datas muda no UC-01 (*a) e a proposta atualizada volta a ser partilhada (passo 1).
  \item[4a.] \textit{Um membro não aceita um item} («dormitório não», \reg{UE05}): a resposta do membro fica visível ao grupo junto desse item, e o sistema mostra ao organizador em que itens há divergência (US-16). O organizador troca a opção (área 2, «Comparar opções») e o grupo vê a proposta atualizada (*a), ou mantém-no e continua no passo 6.
  \item[5a.] \textit{Os membros demoram a responder} (dois dias sem resposta, \reg{UE08; FM.8}): o organizador vê quem ainda falta (US-31) e insiste só com esses. Continua no passo 3 para esses membros.
  \item[5b.] \textit{(Release 2, US-17) O prazo fixado pelo organizador passa e falta alguém responder}: o sistema avisa o organizador de quem falta. Continua em 5a ou no passo 6.
  \item[6a.] \textit{O grupo decide ao vivo, à volta de um só telemóvel} (no café, \reg{UE18; SM.22, PM.29}): o organizador regista o combinado na proposta naquele momento (US-18). Continua no passo 7.
  \item[6b.] \textit{Há membros que ainda não responderam quando o organizador regista a decisão}: a decisão fica registada e esses membros também a veem (US-18). Continua no passo 7.
  \item[*a.] \textit{O organizador altera um item depois de partilhar a proposta}: quem abre a proposta vê a versão atualizada (US-14), com o novo total por pessoa.
  \item[*b.] \textit{O preço de uma opção sobe enquanto o grupo decide} (o voo subiu 19\,€: «malta, subiu», \reg{UE08; FM.2, SM.17}); ponto de extensão \textit{subida de preço}: o sistema avisa o organizador e os membros e mostra de quanto subiu, sem mensagens de urgência (US-08, «Alertar subida de preço»). O total por pessoa é atualizado e o grupo vê-o na proposta (*a).
\end{itemize}

\textbf{Rastreabilidade}
\begin{itemize}
  \item \textbf{Histórias:} US-14 (passos 1 e 2; 2a, *a); US-15 (passo 3; 3a); US-16 (passo 4; 4a); US-31 (passo 5; 3a, 5a); US-18 (passos 6 e 7; 6a, 6b); US-08 (*b); US-17, Release~2 (5b).
  \item \textbf{Sequência consolidada:} atividade 6 «Levar a proposta a quem decide» (intenção: conseguir uma decisão de quem viaja antes que o preço mude) e o seu \textit{breakdown} «o grupo responde devagar e o preço sobe» \reg{U0-09}; atividade 9, «Pôr o grupo a par das contas» \reg{UE17} e «Escolher o resto em conjunto e registá-lo» \reg{UE18}; atividade 7, «Refazer quando uma condição muda» \reg{UE09} (3a).
  \item \textbf{UED, área 5:} «Mostrar ao grupo a ideia e a proposta partilhadas» (passos 1 e 2); «Mostrar ao grupo o que foi cortado e porquê» (2); «Indicar as datas possíveis de cada amigo» (3); «Dar a opinião sobre cada opção» (4); «Registar o que ficou combinado» (6 e 7, 6a); «Avisar o grupo quando o preço sobe» (*b); «Definir um prazo para o grupo decidir» (5b). Ligação da área 5 para a área 1, «para refazer a pesquisa com novas datas» (3a).
  \item \textbf{Afinidade:} «Espero pelo grupo e entretanto o preço sobe»; «A proposta que partilho não chega aos meus amigos».
  \item \textbf{\textit{Breakdowns} (área 5):} FM.2, SM.17 (*b); FM.8 (5a); FM.9, SM.21, AM.25 (2a e passo 2); SM.22, PM.29 (6a).
  \item \textbf{Diagrama de casos de uso (L1):} inclui «Calcular o total por pessoa»; é estendido por «Alertar subida de preço» no ponto de extensão \textit{subida de preço}; ator secundário: membro do grupo.
  \item \textbf{Persona:} Inês.
\end{itemize}
\end{usecase}

\medskip
{\footnotesize\textbf{Notas para revisão.}
(1) O nível é resumo: a decisão demora dias \reg{UE08} e faz-se em várias sessões, de pessoas diferentes (partilhar, responder, insistir, registar).
(2) A extensão 5b depende da US-17, que é da Release~2; fica no caso de uso para mostrar desde já onde entra o prazo.
(3) Na extensão *b, o aviso chega ao organizador e aos membros, como na função do UED «Avisar o grupo quando o preço sobe», e respeita a restrição do UED «Não usar mensagens de urgência para pressionar a decisão», como a US-08.
(4) O ator principal usa o nome das histórias e do diagrama de casos de uso. No modelo de fluxo consolidado, este objetivo corresponde à coordenação do grupo, que nesta entrega é sempre do organizador.
(5) O cenário principal está em colunas porque há três participantes e os membros respondem cada um no seu tempo. O caso em que um membro não pode nas datas propostas e a proposta tem de mudar é um caminho alternativo, por isso está na extensão 3a, que volta ao passo 1.\par}

\end{document}
